\documentclass[conference]{IEEEtran}
\IEEEoverridecommandlockouts
% The preceding line is only needed to identify funding in the first footnote. If that is unneeded, please comment it out.
\usepackage{cite}
\usepackage{amsmath,amssymb,amsfonts}
\usepackage{algorithmic}
\usepackage{graphicx}
\usepackage{textcomp}
\usepackage{xcolor}
\usepackage{booktabs}
\usepackage{multirow}
\usepackage{array}
\renewcommand{\arraystretch}{1.3}

\def\BibTeX{{\rm B\kern-.05em{\sc i\kern-.025em b}\kern-.08em
    T\kern-.1667em\lower.7ex\hbox{E}\kern-.125emX}}
\begin{document}

\title{Language Identification in Low Resource Language of Indian Languages\\
{\footnotesize \textsuperscript{*}Note: Sub-titles are not captured in Xplore and
should not be used}
\thanks{Identify applicable funding agency here. If none, delete this.}
}

\author{\IEEEauthorblockN{1\textsuperscript{st} Nawazish Bilal}
\IEEEauthorblockA{\textit{School of Computer Science and Engineering} \\
\textit{Lovely Professional University}\\
Phagwara, Punjab \\
nawazishkhan378@gmail.com}
\and
\IEEEauthorblockN{2\textsuperscript{nd} Given Name Surname}
\IEEEauthorblockA{\textit{dept. name of organization (of Aff.)} \\
\textit{name of organization (of Aff.)}\\
City, Country \\
email address or ORCID}
\and
\IEEEauthorblockN{3\textsuperscript{rd} Given Name Surname}
\IEEEauthorblockA{\textit{dept. name of organization (of Aff.)} \\
\textit{name of organization (of Aff.)}\\
City, Country \\
email address or ORCID}
\and
\IEEEauthorblockN{4\textsuperscript{th} Given Name Surname}
\IEEEauthorblockA{\textit{dept. name of organization (of Aff.)} \\
\textit{name of organization (of Aff.)}\\
City, Country \\
email address or ORCID}
\and
\IEEEauthorblockN{5\textsuperscript{th} Given Name Surname}
\IEEEauthorblockA{\textit{dept. name of organization (of Aff.)} \\
\textit{name of organization (of Aff.)}\\
City, Country \\
email address or ORCID}
\and
\IEEEauthorblockN{6\textsuperscript{th} Given Name Surname}
\IEEEauthorblockA{\textit{dept. name of organization (of Aff.)} \\
\textit{name of organization (of Aff.)}\\
City, Country \\
email address or ORCID}
}

\maketitle

\begin{abstract}
Efficient parking management is a critical component of modern urban transport systems, directly influencing traffic congestion, fuel consumption, and commuter experience. With the increasing availability of parking occupancy data from smart sensing systems, unsupervised learning techniques can be leveraged to uncover hidden usage patterns without relying on predefined labels. This study presents an unsupervised machine learning approach for analyzing parking occupancy patterns using real-world time-stamped parking data. Feature engineering techniques were applied to extract temporal and occupancy-based attributes, followed by clustering using K-Means and DBSCAN algorithms. The results demonstrate that K-Means effectively identifies dominant parking usage patterns, while DBSCAN excels in detecting irregular and anomalous parking behaviors. A comparative analysis highlights the strengths and limitations of each approach, providing insights for urban transport planning and smart parking management.
\end{abstract}

\begin{IEEEkeywords}
Unsupervised Learning, Parking Occupancy Analysis, K-Means Clustering, DBSCAN, Urban Transport Management
\end{IEEEkeywords}

\section{Introduction}
Urbanization has led to a significant increase in vehicle usage, intensifying the demand for efficient parking management. Poorly managed parking infrastructure contributes to traffic congestion, increased emissions, and wasted commuter time. Traditional parking analysis methods rely heavily on manual surveys or predefined usage categories, which may fail to capture evolving and complex usage patterns.

Unsupervised learning offers a data-driven approach to automatically identify latent structures within parking occupancy data. By clustering parking behavior based on temporal and utilization characteristics, city planners can better understand parking demand dynamics and make informed decisions regarding pricing, infrastructure expansion, and enforcement strategies. This research focuses on analyzing parking occupancy patterns using unsupervised learning techniques applied to real-time parking sensor data. 

\section{LITERATURE REVIEW}

Parking occupancy analysis has gained increasing importance in intelligent transportation systems due to its impact on traffic congestion and urban mobility. Traditional parking studies relied on manual surveys and static models, which lacked scalability and adaptability to dynamic urban environments. With the availability of sensor-based parking data, unsupervised learning techniques have become suitable for discovering latent usage patterns without labeled data.

Clustering methods are widely used for exploratory analysis in such scenarios. K-Means remains one of the most popular centroid-based clustering algorithms due to its simplicity and computational efficiency, though it assumes spherical clusters and is sensitive to noise [1], [6]. Density-based approaches such as DBSCAN, introduced by Ester et al. [2], overcome these limitations by identifying arbitrarily shaped clusters and detecting noise points.

Several studies have applied data-driven techniques to smart parking systems, demonstrating that temporal feature extraction and clustering can effectively reveal parking demand patterns [8], [9]. However, comparative analyses of centroid-based and density-based clustering for parking occupancy analysis remain limited. This study addresses this gap by evaluating K-Means and DBSCAN for urban parking pattern discovery.

\section{METHODOLOGY}
Before you begin to format your paper, first write and save the content as a 
separate text file. Complete all content and organizational editing before 
formatting. Please note sections \ref{AA}--\ref{SCM} below for more information on 
proofreading, spelling and grammar.


Keep your text and graphic files separate until after the text has been 
formatted and styled. Do not number text heads---{\LaTeX} will do that 
for you.

\subsection{Dataset Description}

The dataset used in this study consists of time-stamped parking occupancy records collected from multiple parking zones. The dataset includes the following attributes:
\begin{itemize}
    \item SystemCode: Unique identifier for each parking zone
    \item Capacity: Total number of parking spaces available in the zone
    \item Occupancy: Number of occupied parking spaces at a given time.
    \item Last Updated: Timestamp indicating when the occupancy data was recorded.
\end{itemize}

The dataset does not contain any predefined labels, making it suitable for unsupervised learning approaches.

\begin{table}[htbp]
\caption{Dataset Description}
\centering
\renewcommand{\arraystretch}{1.3}
\begin{tabular}{ll}
\toprule
\textbf{Attribute} & \textbf{Description} \\
\midrule
SystemCode & Unique parking zone identifier \\
Capacity & Total number of parking spaces \\
Occupancy & Number of occupied spaces \\
Last Updated & Timestamp of data collection \\
\bottomrule
\end{tabular}
\label{tab:dataset}
\end{table}

\subsection{Data Preprocessing and Feature Engineering}

Initial exploratory data analysis was performed to inspect data types, identify missing values, and understand feature distributions. The Last Updated attribute was converted to a datetime format to extract meaningful temporal features.

The following features were engineered:
\begin{itemize}
    \item Hour of Day
    \item Day of Week
    \item Weekend Indicator
    \item Occupancy Rate
\end{itemize}
To ensure fair contribution of all features in distance-based clustering algorithms, standardization was applied using z-score normalization.

\begin{table}[htbp]
\caption{Engineered Features Used for Clustering}
\centering
\begin{tabular}{ll}
\toprule
\textbf{Feature} & \textbf{Purpose} \\
\midrule
Occupancy Rate & Normalizes usage across parking zones \\
Hour of Day & Captures temporal parking behavior \\
Day of Week & Identifies weekday trends \\
Weekend Indicator & Separates weekday and weekend usage \\
\bottomrule
\end{tabular}
\label{tab:features}
\end{table}

\subsection{K-Means Clustering}\label{AA}
K-Means clustering was applied to group parking records based on temporal and occupancy behavior. The optimal number of clusters was determined using the Elbow Method, which analyzes the reduction in within-cluster inertia for increasing values of k. Based on the elbow point, three clusters were selected for the final model.

Cluster quality was evaluated using the Silhouette Score, which measures how well data points are separated across clusters.K-Means clustering was applied to group parking records based on temporal and occupancy behavior. The optimal number of clusters was determined using the Elbow Method, which analyzes the reduction in within-cluster inertia for increasing values of k. Based on the elbow point, three clusters were selected for the final model.

Cluster quality was evaluated using the Silhouette Score, which measures how well data points are separated across clusters.

\subsection{DBSCAN Clustering}
DBSCAN (Density-Based Spatial Clustering of Applications with Noise) was employed to identify density-based patterns and detect anomalous parking behavior. Unlike K-Means, DBSCAN does not require prior specification of the number of clusters and can identify noise points.

The eps parameter was selected using a k-distance graph, which revealed a clear elbow point at approximately 0.15. A minimum sample size of five was used to define dense regions.

\subsection{Equations}
Number equations consecutively. To make your 
equations more compact, you may use the solidus (~/~), the exp function, or 
appropriate exponents. Italicize Roman symbols for quantities and variables, 
but not Greek symbols. Use a long dash rather than a hyphen for a minus 
sign. Punctuate equations with commas or periods when they are part of a 
sentence, as in:

% Transformer encoding
\begin{equation}
\mathbf{h}_t = \mathrm{TransformerEncoder}(\mathbf{x}_t)
\end{equation}

% Masked token prediction
\begin{equation}
\hat{\mathbf{y}}_t = \mathrm{softmax}(W_o \mathbf{h}_t + b_o)
\end{equation}

% Masked Language Modeling (MLM) loss
\begin{equation}
\mathcal{L}_{\text{MLM}} = - \sum_{t \in \mathcal{M}}
\log \, P \left( x_t \mid \hat{\mathbf{y}}_t \right)
\end{equation}

% Overall training objective
\begin{equation}
\mathcal{L} = \mathcal{L}_{\text{MLM}}
\end{equation}


\begin{equation}
a+b=\gamma\label{eq}
\end{equation}

Be sure that the 
symbols in your equation have been defined before or immediately following 
the equation. Use ``\eqref{eq}'', not ``Eq.~\eqref{eq}'' or ``equation \eqref{eq}'', except at 
the beginning of a sentence: ``Equation \eqref{eq} is . . .''

\subsection{\LaTeX-Specific Advice}

Please use ``soft'' (e.g., \verb|\eqref{Eq}|) cross references instead
of ``hard'' references (e.g., \verb|(1)|). That will make it possible
to combine sections, add equations, or change the order of figures or
citations without having to go through the file line by line.

Please don't use the \verb|{eqnarray}| equation environment. Use
\verb|{align}| or \verb|{IEEEeqnarray}| instead. The \verb|{eqnarray}|
environment leaves unsightly spaces around relation symbols.

Please note that the \verb|{subequations}| environment in {\LaTeX}
will increment the main equation counter even when there are no
equation numbers displayed. If you forget that, you might write an
article in which the equation numbers skip from (17) to (20), causing
the copy editors to wonder if you've discovered a new method of
counting.

{\BibTeX} does not work by magic. It doesn't get the bibliographic
data from thin air but from .bib files. If you use {\BibTeX} to produce a
bibliography you must send the .bib files. 

{\LaTeX} can't read your mind. If you assign the same label to a
subsubsection and a table, you might find that Table I has been cross
referenced as Table IV-B3. 

{\LaTeX} does not have precognitive abilities. If you put a
\verb|\label| command before the command that updates the counter it's
supposed to be using, the label will pick up the last counter to be
cross referenced instead. In particular, a \verb|\label| command
should not go before the caption of a figure or a table.

Do not use \verb|\nonumber| inside the \verb|{array}| environment. It
will not stop equation numbers inside \verb|{array}| (there won't be
any anyway) and it might stop a wanted equation number in the
surrounding equation.

\subsection{Some Common Mistakes}\label{SCM}
\begin{itemize}
\item The word ``data'' is plural, not singular.
\item The subscript for the permeability of vacuum $\mu_{0}$, and other common scientific constants, is zero with subscript formatting, not a lowercase letter ``o''.
\item In American English, commas, semicolons, periods, question and exclamation marks are located within quotation marks only when a complete thought or name is cited, such as a title or full quotation. When quotation marks are used, instead of a bold or italic typeface, to highlight a word or phrase, punctuation should appear outside of the quotation marks. A parenthetical phrase or statement at the end of a sentence is punctuated outside of the closing parenthesis (like this). (A parenthetical sentence is punctuated within the parentheses.)
\item A graph within a graph is an ``inset'', not an ``insert''. The word alternatively is preferred to the word ``alternately'' (unless you really mean something that alternates).
\item Do not use the word ``essentially'' to mean ``approximately'' or ``effectively''.
\item In your paper title, if the words ``that uses'' can accurately replace the word ``using'', capitalize the ``u''; if not, keep using lower-cased.
\item Be aware of the different meanings of the homophones ``affect'' and ``effect'', ``complement'' and ``compliment'', ``discreet'' and ``discrete'', ``principal'' and ``principle''.
\item Do not confuse ``imply'' and ``infer''.
\item The prefix ``non'' is not a word; it should be joined to the word it modifies, usually without a hyphen.
\item There is no period after the ``et'' in the Latin abbreviation ``et al.''.
\item The abbreviation ``i.e.'' means ``that is'', and the abbreviation ``e.g.'' means ``for example''.
\end{itemize}
An excellent style manual for science writers is \cite{b7}.

\subsection{Authors and Affiliations}
\textbf{The class file is designed for, but not limited to, six authors.} A 
minimum of one author is required for all conference articles. Author names 
should be listed starting from left to right and then moving down to the 
next line. This is the author sequence that will be used in future citations 
and by indexing services. Names should not be listed in columns nor group by 
affiliation. Please keep your affiliations as succinct as possible (for 
example, do not differentiate among departments of the same organization).

\subsection{Identify the Headings}
Headings, or heads, are organizational devices that guide the reader through 
your paper. There are two types: component heads and text heads.

Component heads identify the different components of your paper and are not 
topically subordinate to each other. Examples include Acknowledgments and 
References and, for these, the correct style to use is ``Heading 5''. Use 
``figure caption'' for your Figure captions, and ``table head'' for your 
table title. Run-in heads, such as ``Abstract'', will require you to apply a 
style (in this case, italic) in addition to the style provided by the drop 
down menu to differentiate the head from the text.

Text heads organize the topics on a relational, hierarchical basis. For 
example, the paper title is the primary text head because all subsequent 
material relates and elaborates on this one topic. If there are two or more 
sub-topics, the next level head (uppercase Roman numerals) should be used 
and, conversely, if there are not at least two sub-topics, then no subheads 
should be introduced.

\subsection{Figures and Tables}
\paragraph{Positioning Figures and Tables} Place figures and tables at the top and 
bottom of columns. Avoid placing them in the middle of columns. Large 
figures and tables may span across both columns. Figure captions should be 
below the figures; table heads should appear above the tables. Insert 
figures and tables after they are cited in the text. Use the abbreviation 
``Fig.~\ref{fig}'', even at the beginning of a sentence.

\begin{table}[htbp]
\caption{Table Type Styles}
\begin{center}
\begin{tabular}{|c|c|c|c|}
\hline
\textbf{Table}&\multicolumn{3}{|c|}{\textbf{Table Column Head}} \\
\cline{2-4} 
\textbf{Head} & \textbf{\textit{Table column subhead}}& \textbf{\textit{Subhead}}& \textbf{\textit{Subhead}} \\
\hline
copy& More table copy$^{\mathrm{a}}$& &  \\
\hline
\multicolumn{4}{l}{$^{\mathrm{a}}$Sample of a Table footnote.}
\end{tabular}
\label{tab1}
\end{center}
\end{table}

\begin{figure}[htbp]
\centerline{\includegraphics{fig1.png}}
\caption{Example of a figure caption.}
\label{fig}
\end{figure}

Figure Labels: Use 8 point Times New Roman for Figure labels. Use words 
rather than symbols or abbreviations when writing Figure axis labels to 
avoid confusing the reader. As an example, write the quantity 
``Magnetization'', or ``Magnetization, M'', not just ``M''. If including 
units in the label, present them within parentheses. Do not label axes only 
with units. In the example, write ``Magnetization (A/m)'' or ``Magnetization 
\{A[m(1)]\}'', not just ``A/m''. Do not label axes with a ratio of 
quantities and units. For example, write ``Temperature (K)'', not 
``Temperature/K''.

\section{Results and Discussion}

\subsection{K-Means Results}

K-Means clustering identified three distinct parking usage patterns. One cluster represented high weekday daytime occupancy, typically associated with commercial or business districts. Another cluster showed moderate occupancy during evenings and weekends, indicating mixed-use areas. The third cluster exhibited consistently low occupancy, suggesting underutilized parking zones. The Silhouette Score indicated reasonable cluster separation, demonstrating the effectiveness of K-Means for identifying dominant patterns.

\begin{table}[htbp]
\caption{K-Means Cluster Characteristics}
\centering
\begin{tabular}{cccc}
\toprule
\textbf{Cluster} & \textbf{Avg. Occupancy Rate} & \textbf{Avg. Hour} & \textbf{Interpretation} \\
\midrule
0 & High & Daytime & Commercial / Business Zones \\
1 & Medium & Evening & Mixed-use Areas \\
2 & Low & All-day & Underutilized Zones \\
\bottomrule
\end{tabular}
\label{tab:kmeans}
\end{table}

\subsection{DBSCAN Results}

DBSCAN produced a smaller number of dense clusters along with several noise points. These noise points correspond to irregular parking behavior, such as sudden occupancy spikes during off-peak hours or unusual demand patterns. The Silhouette Score for DBSCAN was observed to be low and slightly negative, which is expected due to overlapping clusters and arbitrary cluster shapes inherent in density-based methods.

\begin{table}[htbp]
\caption{DBSCAN Cluster Distribution}
\centering
\begin{tabular}{cc}
\toprule
\textbf{Cluster Label} & \textbf{Number of Samples} \\
\midrule
Cluster 0 & Normal Dense Region \\
Cluster 1 & Secondary Usage Pattern \\
Noise (-1) & Anomalous Parking Events \\
\bottomrule
\end{tabular}
\label{tab:dbscan}
\end{table}

\subsection{Comparative Analysis}

A comparative evaluation of K-Means and DBSCAN highlights their complementary strengths. K-Means provides compact and interpretable clusters suitable for long-term planning, whereas DBSCAN excels at detecting anomalies and irregular events. While K-Means achieved higher silhouette scores, DBSCAN offered additional insights by isolating exceptional parking behaviors that are critical for real-time monitoring and enforcement.

\begin{table}[htbp]
\caption{Comparison of Clustering Algorithms}
\centering
\begin{tabular}{lll}
\toprule
\textbf{Criterion} & \textbf{K-Means} & \textbf{DBSCAN} \\
\midrule
Cluster Type & Centroid-based & Density-based \\
Requires $k$ & Yes & No \\
Noise Handling & No & Yes \\
Silhouette Score & Higher & Lower \\
Best Use Case & Pattern Discovery & Anomaly Detection \\
\bottomrule
\end{tabular}
\label{tab:comparison}
\end{table}

\subsection{Urban Transport Implications}

The clustering results have practical implications for urban transport management. High-demand clusters can inform dynamic pricing strategies and infrastructure expansion, while underutilized zones may be repurposed or redesigned. Anomalies detected by DBSCAN can assist authorities in identifying event-driven congestion, sensor faults, or enforcement requirements. Integrating both clustering approaches provides a comprehensive understanding of parking demand dynamics.

\section{Conclusion and Future Scope}

This paper presented an unsupervised learning framework for analyzing parking occupancy patterns using K-Means and DBSCAN clustering algorithms. The results demonstrate that K-Means is effective for identifying dominant usage patterns, while DBSCAN is valuable for detecting irregular and anomalous behavior. Future work may incorporate spatial coordinates, apply hierarchical clustering, or integrate predictive models for real-time parking demand forecasting.

\section*{Acknowledgment}

The preferred spelling of the word ``acknowledgment'' in America is without 
an ``e'' after the ``g''. Avoid the stilted expression ``one of us (R. B. 
G.) thanks $\ldots$''. Instead, try ``R. B. G. thanks$\ldots$''. Put sponsor 
acknowledgments in the unnumbered footnote on the first page.

\section*{References}

Please number citations consecutively within brackets \cite{b1}. The 
sentence punctuation follows the bracket \cite{b2}. Refer simply to the reference 
number, as in \cite{b3}---do not use ``Ref. \cite{b3}'' or ``reference \cite{b3}'' except at 
the beginning of a sentence: ``Reference \cite{b3} was the first $\ldots$''

Number footnotes separately in superscripts. Place the actual footnote at 
the bottom of the column in which it was cited. Do not put footnotes in the 
abstract or reference list. Use letters for table footnotes.

Unless there are six authors or more give all authors' names; do not use 
``et al.''. Papers that have not been published, even if they have been 
submitted for publication, should be cited as ``unpublished'' \cite{b4}. Papers 
that have been accepted for publication should be cited as ``in press'' \cite{b5}. 
Capitalize only the first word in a paper title, except for proper nouns and 
element symbols.

For papers published in translation journals, please give the English 
citation first, followed by the original foreign-language citation \cite{b6}.

\begin{thebibliography}{00}
\bibitem{b1} Basu, Joyanta, et al. "Multilingual speech corpus in low-resource eastern and northeastern Indian languages for speaker and language identification." Circuits, Systems, and Signal Processing 40.10 (2021): 4986-5013.
\bibitem{b2} J. Clerk Maxwell, A Treatise on Electricity and Magnetism, 3rd ed., vol. 2. Oxford: Clarendon, 1892, pp.68--73.
\bibitem{b3} I. S. Jacobs and C. P. Bean, ``Fine particles, thin films and exchange anisotropy,'' in Magnetism, vol. III, G. T. Rado and H. Suhl, Eds. New York: Academic, 1963, pp. 271--350.
\bibitem{b4} K. Elissa, ``Title of paper if known,'' unpublished.
\bibitem{b5} R. Nicole, ``Title of paper with only first word capitalized,'' J. Name Stand. Abbrev., in press.
\bibitem{b6} Y. Yorozu, M. Hirano, K. Oka, and Y. Tagawa, ``Electron spectroscopy studies on magneto-optical media and plastic substrate interface,'' IEEE Transl. J. Magn. Japan, vol. 2, pp. 740--741, August 1987 [Digests 9th Annual Conf. Magnetics Japan, p. 301, 1982].
\bibitem{b7} M. Young, The Technical Writer's Handbook. Mill Valley, CA: University Science, 1989.
\end{thebibliography}
\vspace{12pt}
\color{red}
IEEE conference templates contain guidance text for composing and formatting conference papers. Please ensure that all template text is removed from your conference paper prior to submission to the conference. Failure to remove the template text from your paper may result in your paper not being published.

\end{document}
